\chapter{Integrals and Taylor Formulas}

\begin{center}\itshape
$\dfrac\pi2=\dfrac21\cdot\dfrac23\cdot\dfrac43\cdot\dfrac45\cdot\dfrac65\cdot\dfrac67\cdots$ \quad and \quad $\displaystyle\int_0^\infty e^{-x^2}\dd x=\frac{\sqrt\pi}2$.\\
Both come out of one humble family of integrals, which also completes Stirling's formula.
\end{center}

\section{Tools}

\begin{reminder}[integration]
\begin{itemize}
\item \textbf{Integration by parts}: for $u,v$ of class $C^1$ on $[a,b]$, $\int_a^bu'v=[uv]_a^b-\int_a^buv'$.
\item \textbf{Change of variables}: for $\varphi$ of class $C^1$ and $f$ continuous, $\int_{\varphi(a)}^{\varphi(b)}f(x)\dd x=\int_a^bf(\varphi(t))\,\varphi'(t)\dd t$.
\item \textbf{Riemann sums}: for $f$ continuous on $[0,1]$, $\frac1n\sum_{k=1}^nf\left(\frac kn\right)\to\int_0^1f$.
\end{itemize}
\end{reminder}

\begin{reminder}[Taylor formulas]
Let $f$ be of class $C^{n+1}$ on an interval containing $a$ and $x$, and $T_n(x)=\sum_{k=0}^n\frac{f^{(k)}(a)}{k!}(x-a)^k$.
\begin{itemize}
\item \textbf{Integral remainder}: $f(x)=T_n(x)+\displaystyle\int_a^x\frac{(x-t)^n}{n!}f^{(n+1)}(t)\dd t$.
\item \textbf{Taylor--Lagrange inequality}: $|f(x)-T_n(x)|\le\dfrac{|x-a|^{n+1}}{(n+1)!}\sup|f^{(n+1)}|$, the sup over the segment $[a,x]$.
\item \textbf{Taylor--Young}: $f(x)=T_n(x)+o\left((x-a)^n\right)$ (only $C^n$ is needed).
\end{itemize}
\end{reminder}

\section{Wallis integrals}

For $n\in\N$ let $W_n=\displaystyle\int_0^{\pi/2}\sin^nt\dd t$ (equal to $\int_0^{\pi/2}\cos^nt\dd t$ by $t\mapsto\frac\pi2-t$).

\begin{proposition}\label{prop:wallis}
\begin{enumerate}[label=(\roman*)]
\item $W_0=\frac\pi2$, $W_1=1$, and $W_n=\frac{n-1}nW_{n-2}$ for $n\ge2$.
\item $W_{2p}=\dfrac{(2p)!}{4^p(p!)^2}\cdot\dfrac\pi2$ and $W_{2p+1}=\dfrac{4^p(p!)^2}{(2p+1)!}$.
\item $(W_n)$ is decreasing and positive, $nW_nW_{n-1}=\frac\pi2$, and $W_n\sim\sqrt{\dfrac\pi{2n}}$.
\end{enumerate}
\end{proposition}
\begin{proof}
(i) Integrate by parts with $u'=\sin t$, $v=\sin^{n-1}t$:
$W_n=(n-1)\int_0^{\pi/2}\cos^2t\sin^{n-2}t\dd t=(n-1)(W_{n-2}-W_n)$.

(ii) Iterate (i): $W_{2p}=\frac{(2p-1)(2p-3)\cdots1}{(2p)(2p-2)\cdots2}W_0$, and multiply numerator and denominator by $(2p)(2p-2)\cdots2=2^pp!$. Same for odd indices.

(iii) $0\le\sin^{n}t\le\sin^{n-1}t$ on $\left[0,\frac\pi2\right]$, so $(W_n)$ decreases, and $W_n>0$. By (i), $nW_nW_{n-1}=(n-1)W_{n-1}W_{n-2}$, so this quantity is constant, equal to $1\cdot W_1W_0=\frac\pi2$.
Finally $\frac{n-1}n=\frac{W_n}{W_{n-2}}\le\frac{W_n}{W_{n-1}}\le1$, so $W_n\sim W_{n-1}$ and $W_n^2\sim nW_nW_{n-1}/n=\frac\pi{2n}$.
\end{proof}

\begin{theorem}[Stirling's constant]\label{thm:stirlingC}
The constant of Exercise~\ref{ex:stirling} is $C=\sqrt{2\pi}$, so $n!\sim\sqrt{2\pi n}\left(\frac ne\right)^n$.
\end{theorem}
\begin{proof}
Insert $m!\sim C\sqrt m\left(\frac me\right)^m$ into (ii):
\[
W_{2p}=\frac\pi2\cdot\frac{(2p)!}{4^p(p!)^2}\sim\frac\pi2\cdot\frac{C\sqrt{2p}\,(2p)^{2p}e^{-2p}}{4^pC^2p\,p^{2p}e^{-2p}}=\frac\pi2\cdot\frac{\sqrt2}{C\sqrt p}.
\]
By (iii), $W_{2p}\sim\sqrt{\frac\pi{4p}}$. Comparing, $\frac{\pi\sqrt2}{2C}=\frac{\sqrt\pi}2$, so $C=\sqrt{2\pi}$.
\end{proof}

\section{Exercises}

\begin{exercise}[$\star$]\label{ex:xnlnx}
Compute $\displaystyle\int_0^1x^n\ln x\dd x$ for $n\in\N$.
\end{exercise}

\begin{exercise}[$\star$]\label{ex:riemann}
Compute the limits of $\displaystyle\sum_{k=1}^n\frac n{n^2+k^2}$ and of $\displaystyle\left(\prod_{k=1}^n\left(1+\frac kn\right)\right)^{1/n}$.
\end{exercise}

\begin{exercise}[$\star\star$]\label{ex:taylorln}
\begin{enumerate}[label=(\alph*)]
\item Show that $e^x=\sum_{k=0}^\infty\frac{x^k}{k!}$ for every real $x$.
\item Integrating $\frac1{1+t}=\sum_{k=0}^{n-1}(-t)^k+\frac{(-t)^n}{1+t}$ over $[0,1]$, give a second proof that $\sum_{k\ge1}\frac{(-1)^{k+1}}k=\ln2$, with an explicit error bound.
\end{enumerate}
\end{exercise}

\begin{exercise}[$\star\star$]\label{ex:eirr}
Show that $e$ is irrational. \emph{Hint: if $e=\frac pq$, look at $q!\left(e-\sum_{k=0}^q\frac1{k!}\right)$.}
\end{exercise}

\begin{exercise}[$\star\star$]\label{ex:wallisprod}
\emph{Wallis' product.} Show that $\displaystyle\frac{W_{2p}}{W_{2p+1}}\to1$, and deduce
\[
\frac\pi2=\lim_{p\to\infty}\prod_{k=1}^p\frac{(2k)^2}{(2k-1)(2k+1)}=\frac21\cdot\frac23\cdot\frac43\cdot\frac45\cdots
\]
\end{exercise}

\begin{exercise}[$\star\star$]\label{ex:rl}
\emph{Riemann--Lebesgue.} Let $f$ be of class $C^1$ on $[a,b]$. Show that $\displaystyle\int_a^bf(t)\sin(\lambda t)\dd t\to0$ as $\lambda\to\infty$.
\end{exercise}

\begin{exercise}[$\star\star\star$]\label{ex:gaussian}
\emph{The Gaussian integral, without polar coordinates.}
\begin{enumerate}[label=(\alph*)]
\item Show that $\left(1-\frac{x^2}n\right)^n\le e^{-x^2}$ for $0\le x\le\sqrt n$, and $e^{-x^2}\le\left(1+\frac{x^2}n\right)^{-n}$ for $x\ge0$.
\item Show that $\int_0^{\sqrt n}\left(1-\frac{x^2}n\right)^n\dd x=\sqrt n\,W_{2n+1}$ and $\int_0^\infty\left(1+\frac{x^2}n\right)^{-n}\dd x=\sqrt n\,W_{2n-2}$ (for $n\ge1$).
\item Conclude that $\displaystyle\int_0^\infty e^{-x^2}\dd x=\frac{\sqrt\pi}2$.
\end{enumerate}
\end{exercise}

\begin{beauty}
One recurrence, $W_n=\frac{n-1}nW_{n-2}$, gives Wallis' product, the constant in Stirling's formula, and the Gaussian integral. These three
are really the same fact, and in Tome~3 they meet again as the central limit theorem: the $\sqrt{2\pi}$ in the normal density is Stirling's
$\sqrt{2\pi}$. Riemann--Lebesgue (Exercise~\ref{ex:rl}) is the first step of Fourier analysis in Tome~4.
\end{beauty}

\section{Solutions}

\begin{solution}{ex:xnlnx}
On $[\varepsilon,1]$, integrate by parts with $u'=x^n$, $v=\ln x$:
$\int_\varepsilon^1x^n\ln x\dd x=\left[\frac{x^{n+1}}{n+1}\ln x\right]_\varepsilon^1-\int_\varepsilon^1\frac{x^n}{n+1}\dd x$.
As $\varepsilon\to0$, $\varepsilon^{n+1}\ln\varepsilon\to0$, and the integral equals $-\dfrac1{(n+1)^2}$.
\end{solution}

\begin{solution}{ex:riemann}
$\sum_{k=1}^n\frac n{n^2+k^2}=\frac1n\sum_{k=1}^n\frac1{1+(k/n)^2}\to\int_0^1\frac{\mathrm dx}{1+x^2}=\frac\pi4$.
For the product, take logarithms: $\frac1n\sum_{k=1}^n\ln\left(1+\frac kn\right)\to\int_0^1\ln(1+x)\dd x=\left[(1+x)\ln(1+x)-x\right]_0^1=2\ln2-1$, so the limit is $e^{2\ln2-1}=\frac4e$.
\end{solution}

\begin{solution}{ex:taylorln}
(a) Taylor--Lagrange at $a=0$ with $\sup_{[-|x|,|x|]}|\exp^{(n+1)}|=e^{|x|}$: the error is at most $\frac{|x|^{n+1}}{(n+1)!}e^{|x|}\to0$ (the ratio of consecutive terms tends to $0$).

(b) Integrating over $[0,1]$: $\ln2=\sum_{k=0}^{n-1}\frac{(-1)^k}{k+1}+(-1)^n\int_0^1\frac{t^n}{1+t}\dd t$, and
$0\le\int_0^1\frac{t^n}{1+t}\dd t\le\int_0^1t^n\dd t=\frac1{n+1}$. So the partial sums of $\sum\frac{(-1)^{k+1}}k$ are within $\frac1{n+1}$ of $\ln2$.
\end{solution}

\begin{solution}{ex:eirr}
Suppose $e=\frac pq$ with $p,q\ge1$. Then $A=q!\left(e-\sum_{k=0}^q\frac1{k!}\right)$ is an integer, since $q!e=(q-1)!\,p$ and $\frac{q!}{k!}\in\N$ for $k\le q$. But
\[
0<A=\sum_{k>q}\frac{q!}{k!}=\frac1{q+1}+\frac1{(q+1)(q+2)}+\cdots<\sum_{j\ge1}\frac1{(q+1)^j}=\frac1q\le1.
\]
No integer lies strictly between $0$ and $1$.
\end{solution}

\begin{solution}{ex:wallisprod}
By monotonicity, $W_{2p+1}\le W_{2p}\le W_{2p-1}$, so $1\le\frac{W_{2p}}{W_{2p+1}}\le\frac{W_{2p-1}}{W_{2p+1}}=\frac{2p+1}{2p}\to1$. By Proposition~\ref{prop:wallis}(ii),
\[
\frac{W_{2p}}{W_{2p+1}}=\frac\pi2\cdot\frac{(2p)!\,(2p+1)!}{16^p\,(p!)^4}=\frac\pi2\prod_{k=1}^p\frac{(2k-1)(2k+1)}{(2k)^2},
\]
because $(2p)!=\prod_{k=1}^p(2k-1)(2k)$, $(2p+1)!=(2p)!\,(2p+1)$ and $\prod_{k=1}^p(2k)=2^pp!$. Letting $p\to\infty$ gives the product formula.
\end{solution}

\begin{solution}{ex:rl}
Integrate by parts:
$\int_a^bf(t)\sin(\lambda t)\dd t=\left[-f(t)\frac{\cos\lambda t}\lambda\right]_a^b+\frac1\lambda\int_a^bf'(t)\cos(\lambda t)\dd t$.
Both terms are bounded by $\frac1\lambda\left(|f(a)|+|f(b)|+\int_a^b|f'|\right)\to0$.
\end{solution}

\begin{solution}{ex:gaussian}
(a) With $u=\frac{x^2}n$: $\ln(1-u)\le-u$ for $0\le u<1$ gives the first inequality (and it is trivial at $u=1$); $\ln(1+u)\le u$ gives $(1+u)^n\le e^{nu}=e^{x^2}$, the second.

(b) Substitute $x=\sqrt n\sin\theta$, $\theta\in\left[0,\frac\pi2\right]$: $\left(1-\frac{x^2}n\right)^n\dd x=\cos^{2n}\theta\cdot\sqrt n\cos\theta\dd\theta$, giving $\sqrt nW_{2n+1}$.
Substitute $x=\sqrt n\tan\theta$, $\theta\in\left[0,\frac\pi2\right[$: $\left(1+\frac{x^2}n\right)^{-n}\dd x=\cos^{2n}\theta\cdot\frac{\sqrt n}{\cos^2\theta}\dd\theta$, giving $\sqrt nW_{2n-2}$.

(c) Let $I=\int_0^\infty e^{-x^2}\dd x$ (it converges since $e^{-x^2}\le e^{-x}$ for $x\ge1$). By (a) and (b),
$\sqrt nW_{2n+1}\le\int_0^{\sqrt n}e^{-x^2}\dd x\le I\le\sqrt nW_{2n-2}$. By Proposition~\ref{prop:wallis}(iii),
$\sqrt nW_{2n+1}\sim\sqrt n\sqrt{\frac\pi{4n}}=\frac{\sqrt\pi}2$, and likewise for $\sqrt nW_{2n-2}$. Hence $I=\frac{\sqrt\pi}2$.
\end{solution}
